% frontmatter/abstract-th.tex — บทคัดย่อภาษาไทย (มุมวิศวกรรมโยธา)
\chapter*{บทคัดย่อ}
\addcontentsline{toc}{chapter}{บทคัดย่อ}

\noindent\textbf{ชื่อเรื่อง:} \thesistitleTH\\
\textbf{โดย:} \authorname\\
\textbf{อาจารย์ที่ปรึกษา:} \advisor\\

การประมาณราคางานก่อสร้างเป็นกระบวนการสำคัญในการบริหารโครงการก่อสร้าง
โดยเฉพาะในงานภาครัฐที่ต้องอ้างอิงราคากลางตามหลักเกณฑ์ของกรมบัญชีกลาง
และดัชนีราคาวัสดุก่อสร้างของกระทรวงพาณิชย์เพื่อจัดทำบัญชีปริมาณงาน
(Bill of Quantities: BOQ) สำหรับการประกวดราคา การจัดซื้อจัดจ้าง และ
การควบคุมต้นทุนโครงการ. อย่างไรก็ตาม วิศวกรประมาณราคา (Quantity Surveyor)
ในปัจจุบันยังประสบปัญหาสำคัญ 3 ประการ ได้แก่ (1)~ข้อมูลราคาวัสดุกระจายอยู่
ในหลายแหล่งที่ไม่เชื่อมโยงกัน (2)~ราคามาตรฐานภาครัฐมีรอบการปรับปรุง
รายเดือนซึ่งช้ากว่าราคาตลาดจริง และ (3)~ราคาที่นำมาใช้มักเป็นค่าเฉลี่ย
ระดับประเทศที่ไม่สะท้อนความแตกต่างเชิงพื้นที่ของแต่ละจังหวัด.

โครงงานนี้พัฒนาระบบสารสนเทศสำหรับการประมาณราคาวัสดุก่อสร้าง
ที่บูรณาการข้อมูลราคาจาก 11 แหล่งของประเทศไทย ครอบคลุมแหล่งภาครัฐ
3 แห่ง ได้แก่ สำนักงานนโยบายและยุทธศาสตร์การค้า (สนค.) กรมบัญชีกลาง
และกรมการค้าภายใน รวมถึงผู้ค้าปลีกวัสดุก่อสร้างสมัยใหม่ (Modern Trade)
8 แห่ง ได้แก่ HomePro, MegaHome, บุญถาวร, ไทวัสดุ, โกลบอลเฮ้าส์,
ดูโฮม, เอสซีจีโฮม และบีเอ็นบีโฮม. ระบบรองรับการคำนวณต้นทุนต่อหน่วย
ตามหลักเกณฑ์ราคากลางของกรมบัญชีกลาง ครอบคลุมงานก่อสร้างหลัก
5 ประเภท ได้แก่ งานบุกระเบื้องผนัง งานเสาเอ็น/คานทับหลัง งานเหล็กเสริม
งานทาสี และงานก่ออิฐ. ระบบสามารถสร้างรายการวัสดุ (Bill of Materials)
จากปริมาณงานที่ผู้ใช้ระบุ พร้อมแสดงต้นทุนรวมและส่วนประกอบของ
วัสดุแต่ละชนิดเพื่อใช้สนับสนุนการจัดทำเอกสาร BOQ.

ในด้านการจัดการข้อมูลและการประกันคุณภาพราคา ระบบประกอบด้วย
กระบวนการทำความสะอาดข้อมูล การปรับมาตรฐานหน่วยวัด การจับคู่
รายการวัสดุข้ามแหล่งโดยใช้คุณลักษณะมาตรฐาน (ยี่ห้อ ขนาด เกรด)
และการตรวจจับค่าผิดปกติด้วยวิธีทางสถิติ (z-score) ตามแนวทาง
ของ Lee \& Yun (2024) เพื่อกรองราคาที่อาจเกิดจากข้อผิดพลาดทางเทคนิค
หรือโปรโมชันชั่วคราวออกจากการคำนวณ.

การวิเคราะห์ผลเปรียบเทียบราคาวัสดุก่อสร้างใน 6 ภูมิภาค ตลอดช่วงเวลา
9~เดือน พบว่าราคาภาคเอกชนสูงกว่าราคามาตรฐานของกรมบัญชีกลาง
ในกลุ่มปูนซีเมนต์เฉลี่ย 22.9\%, ราคาในภาคใต้สูงกว่าภาคกลาง 22.9\%
ซึ่งสะท้อนต้นทุนการขนส่งระยะไกล, และดัชนีราคาวัสดุก่อสร้าง (CMI)
ของ สนค.\ มีแนวโน้มปรับเพิ่มขึ้นเชิงเส้นเดือนละ 0.4 จุด
($R^{2} = 0.989$) ตลอด 9~เดือนที่ศึกษา. ผลลัพธ์เหล่านี้แสดงให้เห็นว่า
การอาศัยเฉพาะราคามาตรฐานภาครัฐในการประมาณราคามีความเสี่ยง
ที่จะคลาดเคลื่อนจากราคาตลาดจริงอย่างมีนัยสำคัญ ซึ่งอาจกระทบต่อ
ความถูกต้องของงบประมาณและความสำเร็จของการประกวดราคา.

ระบบที่พัฒนาขึ้นจึงเป็นเครื่องมือที่ช่วยวิศวกรโยธา ผู้ประมาณราคา
สถาปนิก และเจ้าหน้าที่ภาครัฐในการตรวจสอบราคาเปรียบเทียบหลายแหล่ง
ได้ภายในเครื่องมือเดียว ลดเวลาในการสืบค้นข้อมูล เพิ่มความแม่นยำ
ของการประมาณการ และยกระดับความโปร่งใสของกระบวนการจัดซื้อจัดจ้าง
ภาครัฐผ่านการเปิดเผยข้อมูลราคาที่สามารถตรวจสอบเปรียบเทียบกันได้.

\noindent\textbf{คำสำคัญ:} การประมาณราคางานก่อสร้าง,
ต้นทุนต่อหน่วยวัสดุ, บัญชีปริมาณงาน (BOQ), ราคากลาง,
ดัชนีราคาวัสดุก่อสร้าง, Modern Trade, การจับคู่รายการวัสดุ
(Canonicalization), การตรวจจับค่าผิดปกติ

\newpage
